\documentclass[10pt,a4paper]{amsart}
\usepackage[margin=19mm]{geometry}
\usepackage{amsmath,amssymb,mathtools}
\usepackage[scr=rsfs,cal=euler]{mathalfa}
\usepackage{xfrac}
\usepackage[unicode,hidelinks,bookmarks=false]{hyperref}
\newcommand{\ie}{i.e.\ }
\newcommand{\cf}{cf.\ }
\newcommand{\resp}{respectively\ }
\newcommand{\Subsection}{\subsection}
\providecommand{\nobreakdash}{\nobreak\hbox{-}}
\setlength{\emergencystretch}{2em}
\multlinegap0pt
\usepackage{amssymb}
\usepackage{dsfont}
\newtheorem{lemma}{Lemma}
\newcommand{\norm}[1]{\lVert#1\rVert}
\title{Multiple mixing and multiple fractional cohomological equation: semisimple setting}
\author{Zhenqi Jenny Wang}
\date{Source: arXiv:2605.21173v1 (CC BY 4.0)}
\begin{document}
\maketitle

\noindent\emph{Source and licence.} Multiple mixing and multiple fractional cohomological equation: semisimple setting. Version arXiv:2605.21173v1 (CC BY 4.0), distributed by arXiv under the Creative Commons Attribution 4.0 International licence, http://creativecommons.org/licenses/by/4.0/, as recorded in the arXiv licence metadata for this exact version and shown on the abstract page https://arxiv.org/abs/2605.21173v1. Full licence notice: \texttt{licenses/arxiv-2605.21173-CC-BY-4.0.txt}.\par\smallskip

\noindent\textit{Editorial scope.} Counted unit: the article's lemma ob:1 (source0.tex lines 2391--2401). The article's complete original proof of it is delivered with it (source0.tex lines 2402--2410, one \texttt{proof} environment of 9 lines). No mathematics is written, repaired or re-derived: the statement and the proof are the article's own lines, in the article's own notation, and no retained line is substituted or re-flowed.  Retained uncounted as the source-local context the proof needs: lines 893--906.  Nothing was omitted from any retained span, so the package carries no omission record.  No retained line contains a citation, so the wrapper carries no bibliography and no bibliography entry.  No retained line contains a figure environment, an image-inclusion command or drawing code, so no image is delivered and the package declares no asset dependency.  Editorial formatting only, in addition: an AMS \texttt{amsart} preamble that restates the article's own declarations and macros used by the retained lines, namely the environments \texttt{lemma} on separate counters, together with the standard packages those lines need.  The retained environments print in this document as Lemma 1 in order of appearance, whereas the article prints the same items with the numbering of its own full document.  The complete native source of the article as distributed in the arXiv e-print for this exact version is bundled unchanged as \texttt{sources/201047/source0.tex}.
\begin{align*}
 W^s(H_\rho)=\int_Z W^s((H_\rho)_z)d\mu(z)
\end{align*}
with respect to the measure $d\mu(z)$.

The existence of the direct integral decompositions allows us to reduce our analysis of the
cohomological equation to irreducible unitary representations. This point of view is
essential for our purposes.

\subsection{Functional calculus for self-adjoint operators}\label{sec:19} Let $H$ be a Hilbert space, and let
  $\mathcal{P}$ be an unbounded, self-adjoint linear operator with a dense domain $\text{Dom}(\mathcal{P})\subset H$. Suppose $\mathcal{P}$ is strictly positive, meaning there is $c>0$ such that
\begin{align*}
  \norm{\mathcal{P}\vartheta}\geq c\norm{\vartheta},\qquad \forall\,\vartheta\in \text{Dom}(\mathcal{P}),
\end{align*}

\begin{lemma}\label{ob:1} Suppose $\mathcal{A}$ is an essentially self-adjoint operator on a Hilbert space $\mathcal{L}$ and satisfies
\begin{align}\label{for:264}
 \langle \mathcal{A}\vartheta, \vartheta\rangle \geq\norm{\vartheta}^2\qquad \forall\, \vartheta\in \text{Dom}(\mathcal{A}).
\end{align} Then
\begin{enumerate}
  \item\label{for:116} if $\mathcal{A}^r\vartheta\in \mathcal{L}$ for some $r>0$ and $\vartheta\in \mathcal{L}$, then $\mathcal{A}^a\vartheta\in \mathcal{L}$ for any $0\leq a\leq r$;

  \smallskip
  \item\label{for:117} for any $r\leq 0$,  $\mathcal{A}^r$ is a bounded linear map on $\mathcal{L}$ with $\norm{\mathcal{A}^r}\leq 1$.
\end{enumerate}
\end{lemma}

\begin{proof} We recall spectral theory in Section \ref{sec:19}. The assumption implies that the spectrum of \(\mathcal{A}\) is contained in \([1,\infty)\). Then for any measurable function \(f\) defined on \([1,\infty)\), one can define the operator \(f(\mathcal{A})\).

\eqref{for:116}:  We note that for any \(0\leq a\leq r\), the function \(x\mapsto x^a\) is dominated by \(x\mapsto x^r\) on \([1,\infty)\).
Then  it follows from the spectral calculus that \(\mathcal{A}^a\vartheta\) belongs to \(\mathcal{L}\).

\eqref{for:117}: If \(r\leq 0\), then for every \(x\geq 1\) we have \(x^r\leq 1\). Consequently, the operator \(\mathcal{A}^r\) is bounded with \(\|\mathcal{A}^r\|\leq 1\).


\end{proof}

\end{document}
